\documentclass[a4paper, serif, 10pt]{moderncv}

%% moderncv
% color and style
\moderncvcolor{black}
\moderncvstyle{banking}

%% page layout/margins
\usepackage[top=2.5cm, bottom=2.5cm, left=1.5cm, right=1.5cm]{geometry}

\nopagenumbers{}

%% Babel config for Italian language
% ref:
% - https://latex3.github.io/babel/guides/locale-italian.html
\usepackage[italian]{babel}

%% fontspec
\usepackage{fontspec}

\IfFontExistsTF{Linux Libertine}{
  \setmainfont[Ligatures={TeX, Common}, Numbers=OldStyle]{Linux Libertine}
}{}
\IfFontExistsTF{Linux Libertine O}{
  \setmainfont[Ligatures={TeX, Common}, Numbers=OldStyle]{Linux Libertine O}
}{}

%% CTeX
% CJK support, used to show CJK characters in the resume
%
% - heading=false: disable ctex heading
% - punct=kaiming: use kaiming punctuations styles for CJK
% - scheme=plain: use plain scheme, do not override \normalsize font size
% - space=auto: space settings for CJK characters
%
% CJK fonts are resolved per locale language, e.g. Noto Serif/Sans CJK JP
% for Japanese resumes, with fallback to the next region if not installed
%
% ref:
% - http://ctan.mirrorcatalogs.com/language/chinese/ctex/ctex.pdf
\usepackage[UTF8, heading=false, punct=kaiming, scheme=plain, space=auto]{ctex}

\IfFontExistsTF{Noto Serif CJK SC}{
  \setCJKmainfont{Noto Serif CJK SC}
}{}
\IfFontExistsTF{Noto Sans CJK SC}{
  \setCJKsansfont{Noto Sans CJK SC}
}{}

%% line spacing
\usepackage{setspace}
\setstretch{1.125}

%% URL styling - use normal text instead of monospace
\urlstyle{same}

%% Patch \httplink and \httpslink to handle full URLs
% The original moderncv macros always prepend http:// or https:// to URLs.
% This causes issues when the URL already has a protocol (e.g., https://example.com)
% resulting in malformed URLs like https://https://example.com.
% This patch checks if the URL already contains "://" and uses it directly if so.
% Note: We use \str_set:Nx (x-expansion) to expand macros like \@homepage first.
\ExplSyntaxOn
\str_new:N \g_yamlresume_url_str

\RenewDocumentCommand{\httpslink}{O{}m}{
  \str_set:Nx \g_yamlresume_url_str {#2}
  \str_if_in:NnTF \g_yamlresume_url_str {://}
    { \url{#2} }
    { \url{https://#2} }
}

\RenewDocumentCommand{\httplink}{O{}m}{
  \str_set:Nx \g_yamlresume_url_str {#2}
  \str_if_in:NnTF \g_yamlresume_url_str {://}
    { \url{#2} }
    { \url{http://#2} }
}
\ExplSyntaxOff

\name{Marco Bianchi}{}
\title{Ingegnere Meccanico Senior specializzato in progettazione e produzione}
\phone[mobile]{+39 02 1234 5678}
\email{marco.bianchi@example.it}
\homepage{https://www.marcobianchi-ingegnere.it}

\address{Via Roma 45, Milano, Lombardia, Italia, 20121}{}{}

\extrainfo{{\small \faLinkedin}\ \href{https://www.linkedin.com/in/marcobianchi-mech}{@marcobianchi-mech}}

\begin{document}

\maketitle

\section{Informazioni personali}

\cvline{}{\begin{itemize}
\item Ingegnere meccanico con oltre 8 anni di esperienza nella progettazione e nello sviluppo di componenti e sistemi meccanici.
\item Specializzato in progettazione CAD 3D, analisi FEM e ottimizzazione di prodotti per i settori automotive e industriale.
\item Esperienza nella gestione di progetti dalla fase di concept fino alla produzione, collaborando con team multifunzionali.
\item Forte orientamento al problem solving, alla qualità e al miglioramento continuo, con conoscenza delle metodologie Lean e Six Sigma.
\end{itemize}}
\section{Formazione}

\cventry{set 2012–lug 2017}
        {Laurea magistrale, Ingegneria Meccanica, Punteggio: 110/110 e lode}
        {Politecnico di Milano}
        {\url{https://www.polimi.it}}
        {}
        {\begin{itemize}
\item Formazione solida in progettazione meccanica, con focus su calcolo strutturale e simulazione numerica.
\item Tesi sperimentale sulla ottimizzazione topologica di componenti per riduzione di peso.
\end{itemize}
\textbf{Corsi}: Progettazione Meccanica Assistita dal Calcolo, Analisi agli Elementi Finiti, Tecnologie di Produzione, Dinamica dei Sistemi Meccanici, Materiali per l'Ingegneria, Gestione della Produzione Industriale}
\section{Esperienza professionale}

\cventry{mar 2020–Presente}
        {Ingegnere Meccanico Senior}
        {Meccanica Avanzata S.p.A.}
        {\url{https://www.meccanicaavanzata.it}}
        {}
        {\begin{itemize}
\item Responsabile della progettazione e dello sviluppo di nuovi componenti meccanici per macchine industriali.
\item Coordinamento di un team di 5 ingegneri e disegnatori, con gestione delle attività di progettazione e validazione.
\item Implementazione di metodologie di analisi FEM per la verifica strutturale e termica dei componenti.
\item Collaborazione con il reparto produzione per ottimizzare i processi di fabbricazione e ridurre i costi.
\item Gestione di progetti chiave, rispettando tempi e budget, con reportistica periodica alla direzione tecnica.
\end{itemize}
\textbf{Parole chiave}: Progettazione meccanica, Analisi FEM, Gestione progetti, Produzione snella, SolidWorks}

\cventry{set 2017–feb 2020}
        {Ingegnere di Progettazione}
        {Ingegneria Dinamica S.r.l.}
        {\url{https://www.ingegneriadinamica.it}}
        {}
        {\begin{itemize}
\item Progettazione di componenti meccanici e assiemi utilizzando CAD 3D (SolidWorks e CATIA).
\item Esecuzione di analisi strutturali e termiche con software FEM per validare i progetti.
\item Redazione di documentazione tecnica, disegni costruttivi e specifiche di prodotto.
\item Supporto al reparto qualità per l'analisi dei difetti e l'implementazione di azioni correttive.
\item Partecipazione a progetti di miglioramento continuo e riduzione dei costi.
\end{itemize}
\textbf{Parole chiave}: CAD 3D, CATIA, Analisi strutturale, Disegno tecnico, Qualità}
\section{Lingue}

\cvline{Italiano}{Madrelingua o bilingue \hfill \textbf{Parole chiave}: Madrelingua}
\cvline{Inglese}{Competenza professionale completa \hfill \textbf{Parole chiave}: Certificazione C1, Inglese tecnico}
\section{Competenze}

\cvline{Progettazione Meccanica}{Esperto \hfill \textbf{Parole chiave}: SolidWorks, CATIA, AutoCAD, Disegno tecnico, Tolleranze}
\cvline{Analisi e Simulazione}{Avanzato \hfill \textbf{Parole chiave}: Analisi FEM, ANSYS, Calcolo strutturale, Termica}
\cvline{Gestione Progetti}{Avanzato \hfill \textbf{Parole chiave}: Pianificazione, Coordinamento team, Metodologie Lean, Six Sigma}
\cvline{Produzione e Materiali}{Intermedio \hfill \textbf{Parole chiave}: Tecnologie di produzione, Metallurgia, Controllo qualità, Manifattura additiva}
\section{Premi}

\cventry{nov 2022}
        {Associazione Italiana di Meccanica}
        {Premio Innovazione Tecnica}
        {}
        {}
        {Riconoscimento per il contributo allo sviluppo di un componente meccanico innovativo che ha ridotto i costi di produzione del 15\%.}
\section{Certificazioni}

\cventry{giu 2019}
        {Dassault Systèmes}
        {Certified SolidWorks Professional}
        {\url{https://www.solidworks.com/certification}}
        {}
        {}

\cventry{ott 2023}
        {Project Management Institute}
        {Project Management Professional (PMP)}
        {\url{https://www.pmi.org/certifications/project-management-pmp}}
        {}
        {}
\section{Pubblicazioni}

\cventry{apr 2021}
        {Ottimizzazione topologica di componenti meccanici per la riduzione di peso}
        {Rivista Italiana di Ingegneria Meccanica}
        {\url{https://www.riim.it/articoli/ottimizzazione-topologica}}
        {}
        {\begin{itemize}
\item Studio sull'applicazione dell'ottimizzazione topologica per la progettazione di componenti leggeri e resistenti.
\item Analisi dei risultati ottenuti tramite simulazioni FEM e validazione sperimentale.
\end{itemize}}
\section{Progetti}

\cventry{gen 2021–dic 2022}
        {Progettazione e prototipazione di un sistema di trasmissione ad alta efficienza per macchine industriali.}
        {Sviluppo di un nuovo sistema di trasmissione}
        {\url{https://www.meccanicaavanzata.it/progetti/trasmissione}}
        {}
        {\begin{itemize}
\item Progettazione meccanica del sistema, inclusi ingranaggi, alberi e supporti.
\item Analisi FEM per la verifica delle sollecitazioni e della durata.
\item Realizzazione di prototipi e test di validazione.
\item Riduzione del 20\% delle perdite energetiche rispetto al sistema precedente.
\end{itemize}
\textbf{Parole chiave}: Trasmissione, Ingranaggi, Prototipazione, Analisi FEM}
\section{Interessi}

\cvline{Motori e Veicoli}{Formula 1, Restauro auto d'epoca, Meccatronica}
\cvline{Tecnologie Emergenti}{Stampa 3D, Materiali compositi, Automazione industriale}
\section{Volontariato}

\cventry{set 2018–Presente}
        {Mentor Tecnico}
        {Associazione Ingegneri Meccanici Italiani}
        {\url{https://www.aimi.it}}
        {}
        {\begin{itemize}
\item Supporto a studenti e giovani ingegneri nello sviluppo di competenze pratiche di progettazione meccanica.
\item Organizzazione di workshop su CAD, FEM e metodologie di progettazione.
\item Partecipazione a eventi di orientamento e divulgazione tecnica.
\end{itemize}}
    
\end{document}